% =====================================================================
%  CHAPTER 8: 食品污染及其预防 (Food Contamination and Prevention)
%  参考：食品主要参考.txt 第1837-2110行 (PRIMARY SOURCE)
% =====================================================================
\chapter{食品污染及其预防}
\chapterintro{
\keyword{掌握}食品污染的概念、来源、分类及对健康的危害；水分活性（Aw）的定义及与微生物生长的关系；常见食品细菌菌属及特征；细菌菌相的概念及食品卫生学意义；菌落总数和大肠菌群的食品卫生学意义及MPN法；黄曲霉毒素的体内代谢、产毒条件、毒性及预防；食品腐败变质的概念、原因、鉴定指标及各类营养素的腐败产物；D值及各种食品保藏方法；$N$-亚硝基化合物、多环芳烃、杂环胺、氯丙醇、丙烯酰胺的污染及预防；食品容器包装材料的卫生问题。\keyword{熟悉}霉菌产毒条件和特点；霉菌菌相；农药和兽药残留。\keyword{了解}霉菌毒素的共同特点；其他霉菌毒素（镰刀菌毒素等）；放射性污染及物理性污染。
}

% =====================================================================
%  第一节 食品污染概述
% =====================================================================
\section{食品污染概述}

\subsection{食品污染的概念与分类}

\begin{keybox}
\textbf{食品污染（food pollution/contamination）：}食品中含有可能对身体健康造成危害或影响其食用价值及商品价值的生物性、化学性及物理性物质，称为食品污染。在食品种植、养殖、生产、加工、储存、运输、销售、烹饪、食用等各个环节中都可能进入或产生于食品中的有害因素，这些过程均可以造成食品污染，而这些有害因素则称为食品污染物（food pollutants）。

\textbf{食品危害的来源（四类）：}
\begin{enumerate}[leftmargin=*]
    \item 食品本身天然含有或自身变化产生的有毒有害物质；
    \item 食品加工储运过程中污染或形成的有害物质；
    \item 外界污染：致病性微生物、寄生虫、农药、重金属等有害化学物及其他污染物污染等；
    \item 食品加工保藏过程中故意加入的成分、食品添加剂、非法添加的有毒有害非食用物质。
\end{enumerate}
\end{keybox}

\subsection{食品污染的分类}

\begin{keybox}
\textbf{食品污染按性质分类（三大类*）：}

\textbf{（1）生物性污染：}微生物、寄生虫和昆虫污染。微生物污染主要是细菌、细菌毒素、霉菌和霉菌毒素以及病毒等的污染，其中细菌和细菌毒素对食品的污染最常见、最重要。病毒污染主要包括甲肝病毒、轮状病毒、诺如病毒和禽流感病毒等。寄生虫和昆虫污染主要有蛔虫、绦虫、囊虫及蝇蛆等。

\textbf{（2）化学性污染：}
\begin{enumerate}[leftmargin=*]
    \item 农药、兽药不合理使用，残留于食品中；
    \item 工业三废、城市生活垃圾未处理随意排放，汞、镉、铅、砷、多环芳烃等有毒有害有机污染物排放到环境中，通过环境转移到食品中；
    \item 食品接触材料、运输工具等接触食品时，食品中的有害物质；
    \item 滥用食品添加剂；
    \item 食品加工储存过程中可能产生的物质，如腌制、烟熏、烘烤等食品过程中产生的$N$-亚硝基化合物、多环芳烃、杂环胺、丙烯酰胺，以及酒中有害的醇类、醛类等；
    \item 掺假、制假过程中加入的物质，如苏丹红、三聚氰胺。
\end{enumerate}

\textbf{（3）物理性污染：}
\begin{enumerate}[leftmargin=*]
    \item 食品的杂物污染：食品生产、加工、储藏、运输、销售等过程中的污染物，包括粮食收割时混入的草籽、液体食品容器中的杂物、食品运输过程中的灰尘等；
    \item 食品的放射性污染：主要来自核爆炸、核废料排放、放射性物质的开采、冶炼、应用及核事故造成的食品污染。转移途径：水—水生生物—植物、动物—人的转移。
\end{enumerate}
\end{keybox}

\subsection{食品污染造成的危害}

\begin{keybox}
\textbf{食品污染造成的危害*：}
\begin{enumerate}[leftmargin=*]
    \item 影响食品的感官性状与营养价值，影响食品生产；
    \item 对身体健康的不良影响，包括急性中毒、慢性危害以及致畸、致突变和致癌作用等。
\end{enumerate}
\end{keybox}

% [补充自往届笔记]
\begin{tipbox}
\textbf{食品安全危害的优先控制顺序（WHO建议）：}微生物污染 $>$ 天然毒素 $>$ 化学性污染 $>$ 食品添加剂。这与公众认知恰好相反——公众往往过分担忧食品添加剂，而忽视微生物污染这一最主要的风险来源。
\end{tipbox}

% =====================================================================
%  第二节 微生物污染及其预防
% =====================================================================
\section{微生物污染及其预防}

\subsection{污染食品的微生物分类}

\begin{keybox}
\textbf{污染食品的微生物按致病能力分类（三大类*）：}

\begin{enumerate}[leftmargin=*]
    \item \textbf{致病性微生物：}能直接对人体致病并造成危害，包括致病性细菌、细菌毒素、人畜共患传染病病原菌和病毒、产毒霉菌和霉菌毒素。
    \item \textbf{条件致病性微生物（相对致病菌）：}通常情况下不致病，但在一定特殊条件下才有效病力的一些微生物。
    \item \textbf{非致病性微生物：}在自然界分布非常广泛，是引起食品腐败变质和食品卫生质量下降的主要原因。
\end{enumerate}
\end{keybox}

\subsection{水分活性（Water Activity, Aw）}

\begin{keybox}
\textbf{水分活性（Water Activity, Aw）的定义*：}在同一条件（温度、湿度、压力等）下，食品水分的蒸汽压（P）与纯水蒸汽压（$P_0$）之比，即 $\text{Aw} = P/P_0$。纯水Aw为1。由于食品中溶质能降低水分的蒸汽压，故Aw值范围为0$\sim$1。Aw反映的是溶液中可被微生物利用的实际含水量，可用来表示食品中可被微生物利用的实际含水量。每个微生物只能在一定的Aw值范围内生长。

\textbf{Aw值与食品中微生物生长繁殖的关系*：}
\begin{enumerate}[leftmargin=*]
    \item 每一种微生物在食品中生长繁殖都有其最低的Aw要求，低于这一要求，微生物的生长繁殖就会受到抑制；
    \item Aw低于0.60时，绝大多数微生物无法生长，Aw小的食品较少出现腐败变质；
    \item 一般说来，细菌生长需要的Aw $>$ 0.9，酵母菌 $>$ 0.87，霉菌 $>$ 0.8；
    \item 同属不同种的微生物对Aw的要求也可不一样；
    \item 细菌形成芽孢时需要比生长时更高的Aw值。
\end{enumerate}
\end{keybox}

% [补充自往届笔记] — supplementary details on Aw
\begin{tipbox}
往届笔记补充：水分活性（water activity, Aw）反映食品中可被微生物利用的水分。每一类微生物在食品中生长繁殖都有其最低的Aw要求，当食品Aw低于这一要求，微生物的生长繁殖就会受到抑制。Aw $<$ 0.6时，绝大多数微生物无法生长，Aw小的食品较少出现腐败变质。细菌 Aw $>$ 0.9，酵母菌 Aw $>$ 0.87，霉菌 Aw $>$ 0.8。
\end{tipbox}

\subsection{常见食品细菌}

\begin{keybox}
\textbf{引起食品腐败变质的常见细菌*：}

\begin{enumerate}[leftmargin=*]
    \item \textbf{假单胞菌属：}需氧性无芽孢杆菌，是冷藏食品（蛋白质变质、冷冻食品）腐败的主要细菌。
    \item \textbf{微球菌属：}为需氧性无芽孢杆菌，是葡萄球菌中常见的腐败菌，常见于未经处理的水和蔬菜中引起腐败。
    \item \textbf{微球菌和葡萄球菌属：}为革兰阳性菌，能产生过氧化氢酶、色素并分解食品中碳水化合物并产生色素。
    \item \textbf{芽孢杆菌属与梭状芽孢杆菌属：}为革兰阳性芽孢杆菌，是罐头食品中常见的腐败菌。
    \item \textbf{肠杆菌科：}为需氧性无芽孢杆菌，与水产制品、肉及蛋的腐败有关。肠杆菌科中除志贺菌属、沙门菌属外，也是常见食品腐败菌。大肠埃希菌是食品中常见的腐败菌，也是食品被粪便污染指示菌之一，变形杆菌属是腐败菌的代表。
\end{enumerate}
\end{keybox}

\subsection{细菌菌相}

\begin{defbox}
\textbf{细菌菌相（Bacterial Flora）：}共存于食品中的细菌种类及其相对数量的构成。其中相对数量较大的细菌称为\textbf{优势菌种（属、株）}。

\textbf{食品卫生学意义：}
\begin{enumerate}[leftmargin=*]
    \item 影响食品腐败变质的特征：使食品的食用价值降低，甚至不能食用。不同优势菌种产生的腐败产物和感官变化各不相同。
    \item 影响食品腐败变质的速度和程度：优势菌种的生长速率和分解能力直接决定腐败进程。在腐败过程中可出现霉菌繁殖，并产生霉菌毒素引起中毒。
    \item 通过分析细菌菌相可判断食品的污染来源和储藏条件。
\end{enumerate}
\end{defbox}

\subsection{菌落总数（Aerobic Plate Count）}

\begin{keybox}
\textbf{菌落总数（Aerobic Plate Count）的定义*：}指单位质量（g）、容积（mL）或表面积（$\text{cm}^2$）的被检样品，在规定条件下（培养基的pH、培养温度与时间、需氧性质等）培养后，所生成的细菌菌落总数。以菌落形成单位（colony forming unit, CFU）表示。

\textbf{食品卫生学意义*：}
\begin{enumerate}[leftmargin=*]
    \item 可作为食品被细菌污染程度及卫生状态的标志；
    \item 是可用于预测食品的存放期限（耐储性）：食品中细菌数量越多，食品腐败变质的速度就越快。
\end{enumerate}
\end{keybox}

\begin{tipbox}
\textbf{注意：}菌落总数低不等于完全无致病菌——因为致病菌可能在低菌落总数条件下存在（如少量沙门菌即可致病）；菌落总数高也不一定代表有致病菌。因此菌落总数是卫生质量的指示指标，而非安全性的直接指标，必须与致病菌检测结合使用。
\end{tipbox}

\subsection{大肠菌群（Coliform Group）与MPN法}

\begin{keybox}
\textbf{大肠菌群（Coliform）的定义*：}指在一定培养条件下能发酵乳糖、产酸产气、需氧和兼性厌氧的革兰阴性无芽孢杆菌。主要包括大肠杆菌科的埃希菌属、柠檬酸杆菌属、肠杆菌属和克雷伯菌属。

\textbf{食品卫生学意义*：}
\begin{enumerate}[leftmargin=*]
    \item \textbf{作为食品受到人与温血动物粪便污染指示菌：}因为大肠菌群直接来源是人与温血动物粪便。典型大肠杆菌的检出说明检品在近期被粪便污染，非典型大肠杆菌的检出说明检品曾经受粪便污染过一段时间（新旧粪便污染）。
    \item \textbf{作为肠道致病菌污染食品的指示菌：}因为大肠菌群与肠道致病菌来源相同，且在一般条件下大肠菌群在外界存活时间与主要肠道致病菌是一致的，食品中大肠杆菌的检出说明食品中有被肠道致病菌污染的可能。
    \item 检测食品加工者、加工设备、存储条件等的污染情况。
\end{enumerate}

\textbf{大肠杆菌检测方法：}
\begin{enumerate}[leftmargin=*]
    \item \textbf{最可能数法（Most Probable Number, MPN）：}当食品中大肠菌群数量较低时采用，相当于每克（mL）食品中大肠菌群的最可能数（MPN）来表示。MPN法是基于泊松分布的一种间接计数方法，是按照一定方法应用统计学原理推算的大肠菌群MPN值。
    \item \textbf{平板计数法：}当食品中大肠菌群数量较高时采用，用平板计数法计大肠菌群的菌落数量表示，表示为每克（mL）检品中大肠菌群的菌落数量CFU/g（mL）。
\end{enumerate}

\textbf{大肠菌群数量单位MPN（most probable number）的定义：}
食品中大肠菌群数量的表示方法，当食品中大肠菌群数量较低时以MPN表示。MPN相当于1g或1mL食品中大肠菌群的最可能数，是按照泊松分布的统计值（GB/T 4789.3-2016），以"MPN/克"或"MPN/毫升"为单位。
\end{keybox}

\subsection{粪便污染指示菌应具备的条件}

\begin{defbox}
\textbf{构成粪便污染指示菌需具备的条件：}
\begin{enumerate}[leftmargin=*]
    \item 仅来自肠道；
    \item 在肠道中数量较多，易于检出；
    \item 在外环境中能生存一定期间；
    \item 食物细菌检验方法敏感简易。
\end{enumerate}

\textbf{肠球菌（粪链球菌）作为粪便污染指示菌的意义：}
肠球菌反映的是低温类食品被粪便污染的指示菌。大肠菌群是常温类食品粪便污染的良好指示菌，但对于低温冷藏/冷冻食品，大肠菌群在低温下易死亡，不宜作为指示菌，而肠球菌对低温抵抗力较强，更适合作为低温类食品的粪便污染指示菌。% [补充自往届笔记]
\end{defbox}

% =====================================================================
%  第三节 霉菌及其毒素对食品的污染
% =====================================================================
\section{霉菌及其毒素对食品的污染}

\subsection{霉菌与霉菌毒素基本概念}

\begin{defbox}
\textbf{霉菌（fungi）：}一类不具叶绿素、无根茎叶分化、具有细胞壁的真核细胞微生物。凡是生长在营养基质上形成绒毛状、蛛网状或絮状菌丝体的真菌，统称为霉菌。

\textbf{霉菌毒素（mycotoxins）：}在霉菌检测中主要是指霉菌污染食品后所产生的有毒代谢产物。

\textbf{霉菌毒素致病特点（与细菌毒素相比）：}
\begin{enumerate}[leftmargin=*]
    \item 毒素较稳定；
    \item 通常通过天然食品中毒；
    \item 中毒以实质器官受损为主；
    \item 没有传染性和免疫性；
    \item 有明显的季节性和地区性特点。
\end{enumerate}
\end{defbox}

\subsection{霉菌产毒特点与产毒条件}

\begin{defbox}
\textbf{霉菌产毒特点*：}
\begin{enumerate}[leftmargin=*]
    \item 霉菌产毒只限于少数的产毒霉菌，而且同一菌种有不同的菌株，产毒能力也不同；
    \item 同一产毒菌株的产毒能力可具有可变性和易变性；
    \item 产毒菌种所产生的霉菌毒素不具有严格的专一性，即一种菌种或菌株可以产生多种不同的毒素，而同一霉菌毒素也可由多种霉菌产生；
    \item 产毒霉菌产生毒素也需要一定的时间。
\end{enumerate}

\textbf{霉菌产毒条件：}
\begin{enumerate}[leftmargin=*]
    \item \textbf{基质：}营养丰富的食品，霉菌生长的可能性大，天然食品上比人工合成培养基更易繁殖。
    \item \textbf{水分：}粮食水分以17\%$\sim$18\%为霉菌生长繁殖的最适条件，粮食Aw值降至0.7以下，一般霉菌均不能生长。
    \item \textbf{湿度：}在不同的相对湿度中，生长繁殖的霉菌种类不同，相对湿度为70\%时霉菌即不能产毒。
    \item \textbf{温度：}霉菌生长繁殖最适宜的温度为25$\sim$30$^\circ$C，在0$^\circ$C以下或30$^\circ$C以上时，霉菌的产毒能力明显减弱或丧失。
    \item \textbf{通风：}大部分霉菌生长繁殖和产毒需要氧气，但毛霉、厌绿曲霉、扩展青霉产毒可在厌氧条件下，能耐受高浓度的CO$_2$。
\end{enumerate}
\end{defbox}

% [补充自往届笔记] — AFB1代谢途径的完整描述
% 往届笔记中有"环氧反应：AFB1-2,3-环氧化物，一部分为解毒过程，一部分发挥毒性、致癌性和致突变性"，这是对代谢途径的重要补充

\subsection{主要产毒霉菌与霉菌毒素}

\begin{defbox}
\textbf{主要产毒霉菌：}
\begin{enumerate}[leftmargin=*]
    \item \textbf{曲霉属：}如黄曲霉、赭曲霉、杂色曲霉；
    \item \textbf{青霉属：}展青霉；
    \item \textbf{镰刀菌属：}禾谷镰刀菌；
    \item \textbf{其他：}黑绿色木霉等。
\end{enumerate}

\textbf{主要霉菌毒素：}黄曲霉毒素、赭曲霉毒素、杂色曲霉毒素、展青霉素、橘青霉素、脱氧雪腐镰刀菌烯醇、玉米赤霉烯酮、伏马菌素、曲酸、单端孢霉烯族化合物、玉米赤霉烯酮、丁烯酸内酯等。
\end{defbox}

\subsection{霉菌污染食品的食品卫生学意义与霉菌菌相}

\begin{defbox}
\textbf{霉菌污染食品的卫生学意义：}
霉菌污染食品的程度以及污染食品的霉菌种类，可从霉菌污染度和霉菌菌相构成两个方面来评定。

\textbf{霉菌菌相（Mold Flora）：}
\begin{enumerate}[leftmargin=*]
    \item \textbf{田野霉}（如气连孢霉、芽枝霉、头孢霉等）占优势：表示粮食新鲜；
    \item \textbf{曲霉和青霉}检出：不表明粮食已霉变，但提示储藏条件可能不良；
    \item \textbf{毛霉、根霉、木霉}检出：说明粮食已开始霉变——这些是储藏后期腐败菌。
\end{enumerate}
\end{defbox}

\subsection{黄曲霉毒素（Aflatoxin, AF）}

\subsubsection{化学结构与理化性质}

\begin{keybox}
\textbf{黄曲霉毒素（Aflatoxin, AF）：}是黄曲霉和寄生曲霉产生的一类有机物。并非所有黄曲霉菌株都能产生AF，但某些寄生曲霉所有菌株几乎全部会产生AF。我国黄曲霉产毒占60\%，寄生曲霉产毒占40\%，是我国长江沿岸和长江以南地区粮食和食品中常见的污染。

\textbf{化学结构：}二呋喃环 + 香豆素（氧杂萘邻酮），基本结构为二呋喃氧杂萘邻酮。AF的毒性与结构有关，含二呋喃环末端有双键者毒性较强并有致癌性。AF的毒性排序：B1 $>$ M1 $>$ G1 $>$ B2 $>$ M2。在天然污染的食品中以AFB1污染最多见，而且其毒性和致癌性也最强，因此在食品卫生监测中常以AFB1作为污染指标。

\textbf{理化性质：}
\begin{enumerate}[leftmargin=*]
    \item \textbf{耐热性强：}一般烹调温度（100$\sim$120$^\circ$C）几乎不被破坏，需\textbf{268$\sim$269$^\circ$C}才完全分解；
    \item \textbf{脂溶性：}难溶于水（溶解度仅10$\sim$20 mg/L），易溶于油脂和多种有机溶剂（氯仿、甲醇等）；
    \item \textbf{在中性或弱酸性溶液中较稳定：}在强碱性条件下（pH 9$\sim$10），内酯环被打开，可被破坏。此原理被用于碱炼去毒；
    \item \textbf{紫外线敏感：}在紫外线下产生荧光（B族发蓝紫色荧光，G族发黄绿色荧光），可用于检测。
\end{enumerate}
\end{keybox}

\subsubsection{产毒条件与污染食品}

\begin{keybox}
\textbf{产毒条件*：}
\begin{itemize}[leftmargin=*]
    \item \textbf{温度：}黄曲霉生长繁殖的温度范围12$\sim$42$^\circ$C，最适产毒温度为25$\sim$33$^\circ$C；
    \item \textbf{水分：}最适Aw值为0.93$\sim$0.98；
    \item \textbf{基质：}花生、玉米（产毒量最高），其次为大米、小麦、大麦、豆类等。
\end{itemize}

\textbf{易污染食品：}
\begin{itemize}[leftmargin=*]
    \item AF主要污染粮油及其制品，其中以花生和玉米污染最为严重，其次为棉籽油，再次为大米、小麦、大麦、豆类等；
    \item 干果类食品（核桃、杏仁、榛子）；
    \item 动物性食品（乳及乳制品、肝脏、干咸鱼）；
    \item 家庭自制发酵食品；
    \item 奶牛食入被AFB1污染的饲料后，牛乳中排出AFB1的羟基化代谢产物AFM1。牛乳及乳制品中的AFM1比较稳定，且目前一般的家庭加热及加工过程均对AFM1水平无显著影响。
\end{itemize}
\end{keybox}

\subsubsection{黄曲霉毒素的体内代谢（重点掌握）*}

\begin{keybox}
\textbf{AFB1的体内代谢过程*\difficulty：}

AFB1进入机体后主要在肝脏进行代谢，在肝细胞微粒体混合功能氧化酶系（CYP450）催化下，发生羟化、脱甲基和环氧化反应，其中主要的羟基化产物为终致癌物。

\textbf{代谢途径：}
\begin{enumerate}[leftmargin=*]
    \item \textbf{AFB1 $\rightarrow$ AFM1（羟化反应）：}AFB1在肝微粒体酶催化下的羟化产物，存在于牛乳和乳制品中。奶牛摄入AFB1污染饲料后，乳汁中排出AFM1。
    \item \textbf{AFB1 $\rightarrow$ AFQ1（羟化反应）：}AFB1的另一种羟化产物，将其转化为AFQ1是一种解毒过程。
    \item \textbf{AFB1 $\rightarrow$ AFH1（还原反应）：}AFQ1的还原产物，黄曲霉毒素醇转化为AFH1可能也是一种解毒过程。
    \item \textbf{AFB1 $\rightarrow$ AFB1-8,9-环氧化物（环氧化反应）：}由CYP450催化的环氧化反应生成，是AFB1的终致癌物。可与DNA结合形成AFB1-N$^7$-鸟嘌呤加合物，导致DNA损伤和突变。% [补充自往届笔记] AFB1-2,3-环氧化物，一部分为解毒过程，一部分发挥毒性、致癌性和致突变性
\end{enumerate}

AF有很强的急性毒性，也有明显的慢性毒性与致癌性。AF对肝脏有特殊亲和性（嗜肝毒），具有较强的肝脏毒性并有致癌作用。
\end{keybox}

\subsubsection{黄曲霉毒素的毒性}

\begin{keybox}
\textbf{黄曲霉毒素的毒性表现（三方面*）：}
\begin{enumerate}[leftmargin=*]
    \item \textbf{急性毒性：}
    \begin{itemize}[leftmargin=*]
        \item 主要靶器官：\textbf{肝脏}——引起肝细胞坏死、胆管增生、肝出血（急性中毒性肝炎）。
        \item 中毒临床表现以黄疸为主，出现发热、呕吐和食欲缺乏，严重者出现腹水、下肢水肿、肝脾大及肝硬化。
        \item AFB1经口LD$_{50}$：鸭雏 0.24$\sim$0.36 mg/kg（最敏感物种之一），大鼠 5.5$\sim$17.9 mg/kg。
    \end{itemize}
    \item \textbf{慢性毒性：}
    \begin{itemize}[leftmargin=*]
        \item 主要表现为生长发育障碍，肝脏出现亚急性或慢性损伤，肝功能降低，肝实质细胞变性、坏死，形成再生结节。
        \item 生长发育迟缓，影响蛋白质合成和营养利用。
        \item 免疫功能下降，抑制细胞免疫和体液免疫。
    \end{itemize}
    \item \textbf{致癌性（已知最强的天然致癌物！）：}
    \begin{itemize}[leftmargin=*]
        \item 主要是肝癌，诱发肝细胞型肝癌。
        \item 诱癌强度是二甲基亚硝胺的\textbf{75倍}。
        \item 可诱发多种实验动物（大鼠、小鼠、鸭、鱼、猴等）的多器官癌症，无动物种属对AF致癌完全抵抗。
        \item 灵长类动物亦可诱癌——猴经AFB1处理可发生肝癌。
        \item \textbf{与人类肝癌密切相关：}膳食AF暴露量与原发性肝癌发病率呈正相关。\textbf{尤为重要的是AF与HBV（乙肝病毒）的协同致癌作用！}HBV感染者暴露于AF时，肝癌风险急剧增高。
        \item IARC分类：AFB1为\textbf{Group 1（确认人类致癌物）}。
    \end{itemize}
\end{enumerate}
\end{keybox}

\subsubsection{黄曲霉毒素的预防措施}

\begin{keybox}
\textbf{黄曲霉毒素的预防措施*：}

\textbf{（1）食物防霉（最根本的措施）：}控制粮食水分含量在安全水分以下，注意低温保藏，注意通风。

\textbf{（2）去除毒素（7种方法）：}
\begin{enumerate}[leftmargin=*]
    \item \textbf{挑选霉粒法：}对花生、玉米去毒效果好；
    \item \textbf{碾轧加工法：}将受污染的大米加工成精米，可降低毒素含量；
    \item \textbf{加水搓洗法：}反复搓洗去除表面毒素；
    \item \textbf{植物油加碱去毒：}AF与NaOH反应，内酯结构中的酯键被水解破坏，形成香豆素钠盐。将碱炼后的水洗可去除毒素。本方法用于植物油去毒效果最好。
    \item \textbf{物理吸附法：}含毒素的植物油加入活性白陶土或活性炭等吸附剂，适用于液体食品（如植物油），效果较好，对固体食品效果不明显。
    \item \textbf{氨气处理法：}在18kg压力、72$\sim$82$^\circ$C状态下，氨气处理可使AF被去除98\%$\sim$100\%，但可使粮食中蛋白质含量增加，同时可能破坏维生素。
    \item \textbf{日光或紫外线照射法：}破坏AF的化学结构。
\end{enumerate}

\textbf{（3）制定食品中AF限量标准：}具体限量标准执行GB 2761。
\end{keybox}

\subsection{镰刀菌毒素}

\begin{defbox}
\textbf{镰刀菌毒素（Fusarium Toxins）：}由镰刀菌属（Fusarium spp.）产生的有毒次级代谢产物，根据其化学结构可分为单端孢霉烯族化合物、玉米赤霉烯酮、伏马菌素和丁烯酸内酯等。

\textbf{1. 单端孢霉烯族化合物：}
\begin{itemize}[leftmargin=*]
    \item 主要有T-2毒素、二乙酸藨草镰刀菌烯醇、雪腐镰刀菌烯醇和脱氧雪腐镰刀菌烯醇。T-2毒素是三线镰刀菌和拟枝孢镰刀菌的代谢产物，是食物中毒性白细胞缺乏症（ATA）的病原物质。
    \item 化学性质：化学性质稳定，在通常条件下不产生荧光，耐热，在中性和酸性环境中不被破坏。
    \item 毒性特点与食品卫生学意义：具有强细胞毒性、免疫抑制及致畸作用，对人体有潜在的致癌性，急性毒性强，可引起动物呕吐。主要污染麦类、玉米及麦类制品。
\end{itemize}

\textbf{2. 玉米赤霉烯酮（Zearalenone, ZEN）：}
\begin{itemize}[leftmargin=*]
    \item \textbf{来源：}又称F-2毒素，由禾谷镰刀菌、黄色镰刀菌、雪腐镰刀菌等产生。
    \item \textbf{毒性特点：}具有雌激素样作用，可表现出对生殖系统的毒性作用。对动物主要污染玉米，其次为小麦、大麦、大米等粮食作物。
\end{itemize}

\textbf{3. 伏马菌素（Fumonisins）：}
\begin{itemize}[leftmargin=*]
    \item \textbf{来源：}主要由串珠镰刀菌产生，可分为伏马菌素B1（FB1）和伏马菌素B2（FB2）两大类。食品中以FB1污染为主，主要污染玉米及玉米制品。FB1通过干扰神经鞘磷脂和鞘氨醇的合成，影响鞘脂类代谢。
    \item \textbf{主要危害：}具有神经毒性作用，可引起马的脑白质软化症；对多种动物具有肾脏毒性；伏马菌素是促癌物，主要诱发原发性肝癌。
\end{itemize}
\end{defbox}

\subsection{其他重要霉菌毒素}

\begin{defbox}
主要霉菌毒素及其毒性（了解）：
\begin{itemize}[leftmargin=*]
    \item \textbf{赭曲霉毒素（Ochratoxin A, OTA）：}主要由赭曲霉和纯绿青霉产生，主要靶器官为肾脏，可引起巴尔干肾病。IARC Group 2B。
    \item \textbf{展青霉素（Patulin）：}扩展青霉产生，污染水果及制品（苹果汁、苹果酱最突出），具有神经毒性、致突变性和免疫抑制性。
    \item \textbf{杂色曲霉毒素：}可引起肝癌等。
\end{itemize}
\end{defbox}

% =====================================================================
%  第四节 食品腐败变质
% =====================================================================
\section{食品腐败变质}

\subsection{食品腐败变质的概念与原因}

\begin{keybox}
\textbf{食品腐败变质（food spoilage）：}指食品在以微生物为主的各种环境因素作用下，食品的成分被分解、破坏，失去或降低食用价值的一切变化。包括肉、鱼、禽、蛋的腐败，粮食的霉变、蔬菜水果的溃烂、油脂酸败、牛奶的变质等。

\textbf{食品腐败变质的原因（三类因素*）：}

\textbf{（1）微生物因素：}主要是细菌、霉菌和酵母菌。
\begin{itemize}[leftmargin=*]
    \item 分解蛋白质的微生物：主要是细菌，其次是霉菌和酵母菌；
    \item 分解碳水化合物的微生物：主要有细菌和酵母菌，能分解某些碳水化合物的细菌；
    \item 分解脂肪的微生物：有能产生脂肪酶，能分解脂肪的革兰阴性菌，食品中常见的有霉菌（如黄曲霉、腊叶芽枝霉等）和假丝酵母菌等；其中酵母菌主要是解脂假丝酵母，对糖类不发酵，分解脂肪和蛋白质的能力却很强。
\end{itemize}

\textbf{（2）食物本身的组成和性质：}
\begin{itemize}[leftmargin=*]
    \item 食品中的酶：食品经酶作用，引起食品组成成分的分解，加速食品的腐败变质；
    \item 食品的营养成分和水分：食品中含有丰富的营养成分，是微生物的良好培养基，富含蛋白质的动物性食品主要以蛋白质腐败为基本特征；碳水化合物含量高的食品，主要在细菌和酵母菌的作用下，以产酸发酵为基本特征；食品中Aw值越小，微生物越不易繁殖，食品越不易腐败变质；
    \item 食品的物理性质：食品pH高低也是制约微生物生长影响食品腐败变质的重要因素之一，一般在酸性或碱性环境中都不利于细菌生长。
\end{itemize}

\textbf{（3）环境因素：}温度、湿度、氧气、光照（紫外线）等。
\end{keybox}

\subsection{食品腐败变质过程中营养成分的变化*}

\begin{keybox}
\textbf{（1）食品中蛋白质的分解：}食品中的蛋白质在微生物的蛋白酶、肽链内切酶等作用下可分解形成氨基酸，氨基酸及胺等在低分子量下可通过脱羧基、脱氨基等作用，形成各种腐败产物。
\begin{itemize}[leftmargin=*]
    \item 细菌脱羧酶作用下：酪氨酸、组氨酸、赖氨酸可分别生成酪胺、组胺和尸胺，具有恶臭味。
    \item 微生物脱氨基酶作用下：生成丁酸、己酸等。
    \item 色氨酸脱羧基色胺、甲基色胺，具有粪臭味。
    \item 含硫氨基酸在脱硫醇酶作用下生成具有恶臭味的硫醇。
\end{itemize}

\textbf{（2）食品中脂肪的酸败（rancidity）：}指食品中脂肪的腐败变质现象。脂肪腐败的化学过程复杂，主要是经水解和氧化反应产生的分解产物。
\begin{itemize}[leftmargin=*]
    \item 中性脂肪分解为甘油和脂肪酸，后者可进一步氧化为低级脂肪酸、醛、酮等；
    \item 不饱和脂肪酸中双键被氧化形成过氧化物，进一步分解为醛、酮、酸，含不饱和脂肪酸越高的食品越易氧化。
    \item 酸败过程中形成的醛、酮和某些羧酸使酸败的脂肪带有特殊的刺激性臭味，即所谓的"哈喇"味（油脂气味）。含不饱和脂肪酸越高的食品越易被氧化。
    \item 脂肪分解的早期，酸败尚不明显，即出现酸度增高，脂肪酸酸度与醛酮含量的反应很敏感。
\end{itemize}

\textbf{酸败的鉴定：}
\begin{enumerate}[leftmargin=*]
    \item 油脂气味；
    \item 过氧化值上升——脂肪酸败的最早期指标；
    \item 酸价和羰基价反应阳性；
    \item 碘、醇等食用脂肪氧化。
\end{enumerate}

\textbf{（3）碳水化合物的分解：}碳水化合物在微生物和植物组织酶的作用下，可分解成单糖、醇、醛、酮、CO$_2$和水等。该过程的主要变化是食品酸度升高，产生酸、醇、醛等气味。
\end{keybox}

\subsection{食品腐败变质的鉴定指标}

\begin{keybox}
\textbf{食品腐败变质的鉴定指标（四类*）：}

\textbf{（1）感官鉴定（最敏感的方法）：}色、香、味、形，主要通过视觉、嗅觉、味觉和触觉来鉴定食品腐败变质。例如：食品腐败常出现不良气味、颜色出现褐色、绿色或失去光泽，组织软化或发黏等。

\textbf{（2）物理指标：}常用的有油脂比重、折光率、黏度等。

\textbf{（3）化学鉴定：}通过直接测定微生物代谢产物的腐败产物。
\begin{itemize}[leftmargin=*]
    \item \textbf{挥发性盐基总氮（Total Volatile Basic Nitrogen, TVBN）：}指蛋白质丰富食品的水浸液在弱碱性条件下与水蒸气一起蒸馏出来的总氮量，在此过程中蒸馏出来的氨和形成的胺的含氮物。TVBN与食品腐败变质程度之间有明确的对应关系，为鱼、肉类蛋白腐败鉴定的化学指标。
    \item \textbf{三甲胺：}主要用于测定鱼、虾等水产品的新鲜程度。
    \item \textbf{组胺：}细菌脱羧酶作用于组氨酸生成组胺，当鱼中组胺超过200 mg/100g就可引起人类过敏性食物中毒。
    \item \textbf{K值（K value）：}指ATP分解的低级产物肌苷（HxR）和次黄嘌呤（Hx）占ATP系列分解产物ATP$+$ADP$+$AMP$+$IMP$+$HxR$+$Hx的百分比。主要用于鉴定鱼类早期腐败。$K < 20\%$说明鱼体新鲜，$K \geq 40\%$说明鱼体开始有腐败迹象。
    \item \textbf{pH值：}一般腐败开始时微生物生长受限，氨基酸脱羧后生成胺类物质，pH值出现V字形变动。
    \item \textbf{过氧化值与酸价：}过氧化值是脂肪酸败鉴定早期指标，酸价反映脂肪酸败的程度。
\end{itemize}

\textbf{（4）微生物检验：}常用指标为菌落总数、大肠菌群和细菌菌相。
\end{keybox}

\subsection{食品腐败变质的食品卫生学意义及处理原则}

\begin{keybox}
\textbf{食品卫生学意义：}
食品腐败变质可使感官性状发生改变，造成食品营养成分分解，营养价值严重降低；进而，腐败变质食品必然受到致病微生物等严重污染，并且可能带有致病菌和产毒霉菌，食用后易引起不良反应和食物中毒。

\textbf{腐败变质食品的处理原则：}
\begin{enumerate}[leftmargin=*]
    \item \textbf{严重腐败变质：}销毁或改作工业用；
    \item \textbf{轻度腐败变质的食品：}经适当的加工处理，将变质的主要指标降至安全限量以下，可以限期食用；
    \item \textbf{已经变质的食品：}禁止食用；
    \item \textbf{局部变质食品：}挑选去除变质部分，利用其他完好部位。
\end{enumerate}
\end{keybox}

% =====================================================================
%  第五节 食品保藏方法
% =====================================================================
\section{食品保藏方法}

\begin{defbox}
\textbf{食品保藏（Food Preservation）：}指为了防止食品腐败变质，延长食品可食用的期限，对食品中的微生物进行杀灭或抑制其生长繁殖，并对食品进行加工处理，通过改变食品的温度、水分、氢离子浓度、渗透压以及采用其他抑菌杀菌措施，以达到防止食品腐败变质的目的。
\end{defbox}

\subsection{化学保藏法}

\begin{keybox}
\textbf{化学保藏法：}
\begin{enumerate}[leftmargin=*]
    \item \textbf{盐腌法：}利用高渗透压抑制微生物生长，食盐浓度达10\%即可抑制大多数细菌，食品含糖量达到60\%$\sim$65\%即可。
    \item \textbf{酸渍法：}多数微生物在pH 4.5以下不能生长繁殖，故可利用控制氢离子浓度来抑制微生物生长，多用于各种蔬菜，如泡菜和酸菜等。
    \item \textbf{防腐剂：}在食品中可添加符合食品添加剂标准的防腐剂，来抑制或杀灭食品中引起腐败变质的微生物，如苯甲酸、山梨酸等。抗氧化剂用于防止油脂酸败。
\end{enumerate}
\end{keybox}

\subsection{低温保藏法}

\begin{keybox}
\textbf{低温保藏法：}
\begin{itemize}[leftmargin=*]
    \item \textbf{冷藏：}指在不冻结状态下的低温储藏，温度一般设定在$-1\sim10^\circ$C范围内。其原理是低温（10$^\circ$C以下）可抑制微生物的繁殖，降低食品中酶的活性，可有效延缓食品的变质。
    \item \textbf{冷冻储藏：}指在$-18^\circ$C以下保藏。在此温度下几乎所有的微生物不再繁殖，冷冻食品可以较长期地保藏。在工艺上，速冻（急冻）、缓慢解冻（缓化）有利于保持食品的品质。
\end{itemize}

\textbf{冰晶生成带：}
\begin{itemize}[leftmargin=*]
    \item $-1\sim-5^\circ$C为冰晶生成带，可导致食品中水结冰率为85\%，不利于保持食品品质。
    \item $-8\sim-120^\circ$C为冻结带，食品中水结冰率为90\%。速冻（急冻），迅速通过冰晶生成带，形成的冰晶数量多、体积小而不损伤食品组织细胞。
\end{itemize}
\end{keybox}

\subsection{高温杀菌保藏法}

\begin{keybox}
\textbf{高温杀菌保藏法：}

\textbf{D值（Decimal Reduction Time）*\difficulty：}在一定温度下，能杀死食品中某种细菌90\%的时间（min）。主要用于比较细菌的死亡速度。如金黄色葡萄球菌 D$_{63}$ = 7；肉毒梭菌芽孢 D$_{121}$ = 0.1$\sim$0.23。

\textbf{常用高温杀菌方法：}
\begin{enumerate}[leftmargin=*]
    \item \textbf{常压杀菌法：}温度控制在100$^\circ$C及以下，杀灭致病菌和繁殖型微生物，适用于液态食物。常用巴氏杀菌法处理食品（如牛乳、啤酒、酱油、葡萄酒等）。以牛乳为例，低温巴氏杀菌法温度63$^\circ$C，30分钟；高温短时巴氏杀菌法温度72$^\circ$C，15秒。
    \item \textbf{高压杀菌法：}温度100$\sim$121$^\circ$C（绝对压力为0.2 MPa），适用于肉类制品、罐头食品，可杀灭繁殖型和芽孢型细菌。
    \item \textbf{超高温瞬时杀菌（UHT）：}在封闭的系统中加热到120$^\circ$C以上，持续几秒钟后迅速冷却至常温的一种杀菌方法，用于牛乳。
    \item \textbf{微波杀菌：}目前多用915 MHz和2450 MHz两个频率。915 MHz可获得较大穿透深度，适用于加热含水量高、厚度或体积较大的食品；2450 MHz适用于含水量低的食品。

    特点：快速、节能、对食品品质影响小，能保留更多活性物质和营养成分。
\end{enumerate}
\end{keybox}

\subsection{食品的干燥脱水保藏}

\begin{keybox}
\textbf{食品脱水保藏的机制：}降低食品水分至15\%以下或Aw值在0.00$\sim$0.60之间，可抑制各种微生物生长，使食品在常温下长期保藏。

\textbf{常用方法：}日晒、阴干、喷雾干燥、减压蒸发、冷冻干燥等。

\textbf{酶性褐变：}是酚酶催化酚类物质形成醌及其聚合物的结果。

\textbf{非酶褐变（美拉德反应/Maillard反应）：}由蛋白质、肽和氨基酸等的氨基与糖及脂肪氧化的醛、酮等羰基所发生的反应，使食物呈现褐色和焦香味。影响因素：加热与给氧促进褐变；水分活性（Aw 0.6最适于褐变）；SO$_2$可与羰基结合而阻止褐变。
\end{keybox}

\subsection{食品辐照保藏}

\begin{keybox}
\textbf{辐照保藏：}

\textbf{辐照灭菌（radappertization）：}指使用波长小于200 nm的电磁波辐照处理食物，达到商业无菌要求的一种处理方式。用高剂量（10$\sim$15/50 kGy）杀灭食品中所有微生物。辐照可使食品中的微生物数量减少到零或有限个数，且不会因辐照而产生放射性。

\textbf{辐照消毒（radicidation）：}使用适当的中等剂量（5$\sim$10 kGy），杀死无芽孢致病菌，主要要求杀灭沙门菌以及各种微生物，辐照后的食品3天内不会发虫、霉变。

\textbf{辐照防腐（radurization）：}使用剂量1$\sim$5 kGy，主要杀灭腐败菌，延长新鲜食品的后熟期及保藏期。主要用于抑制发芽、延长后熟期、杀虫等。
\end{keybox}

% =====================================================================
%  第六节 农药和兽药残留
% =====================================================================
\section{农药和兽药残留}

\subsection{农药的定义与分类}

\begin{keybox}
\textbf{农药（pesticides）的定义：}指用于预防、控制危害农业、林业的病、虫、草、鼠和其他有害生物以及有目的地调节植物、昆虫生长的化学合成或者来源于生物、其他天然物质的一种物质或者几种物质的混合物及其制剂。

\textbf{农药分类（按五种方式）：}
\begin{enumerate}[leftmargin=*]
    \item \textbf{按用途分：}杀虫剂、杀鼠剂、杀菌剂、杀线虫剂、杀螺剂、除草剂、植物生长调节剂等。
    \item \textbf{按化学组成及结构分：}有机磷类、有机氯类、氨基甲酸酯类、拟除虫菊酯类、有机砷类、有机汞类等。
    \item \textbf{按成分和来源分：}无机类、有机类和生物农药。
    \item \textbf{按急性毒性大小分：}剧毒、高毒、中等毒和低毒农药。
    \item \textbf{按残留特性分：}高残留、中等残留和低残留农药。
\end{enumerate}
\end{keybox}

\subsection{兽药的定义}

\begin{defbox}
\textbf{兽药（veterinary drugs）：}指用于预防、治疗、诊断动物疾病或者有目的地调节动物生理机能的物质，包括药物饲料添加剂。主要包括血清制品、疫苗、诊断制品、微生态制品等。
\end{defbox}

\subsection{农药残留与最大残留限量}

\begin{keybox}
\textbf{农药残留（pesticide residues）：}指任何由于使用农药而在食品、农产品和动物饲料中出现的特定物质，包括被认为具有毒理学意义的农药衍生物，如农药转化物、代谢物、反应产物及杂质等。

\textbf{最大残留限量（Maximum Residue Limits, MRLs）：}指在食品或农产品内部或表面，按照良好农业生产规范（GAP）使用农药后，允许农药残留的最高浓度，以每千克食品或农产品中农药残留的毫克数（mg/kg）表示。
\end{keybox}

\subsection{农药污染食品的途径与预防}

\begin{defbox}
\textbf{农药污染食品的途径：}
\begin{enumerate}[leftmargin=*]
    \item 直接污染（施药）——最主要的污染途径；
    \item 农作物从污染环境中吸收；
    \item 通过食物链和生物富集作用污染食品；
    \item 储运过程中的污染；
    \item 其他来源的污染——食品加工、储藏、运输、销售过程中的污染。
\end{enumerate}

\textbf{预防控制措施：}
\begin{enumerate}[leftmargin=*]
    \item 加强农药管理；
    \item 安全合理使用农药；
    \item 制定和执行食品中农药残留限量标准；
    \item 研制和推广高效低毒低残留新品种。
\end{enumerate}
\end{defbox}

% [补充自往届笔记] — 农药分类毒性详解
\begin{tipbox}
\textbf{往届笔记补充——各类农药毒性特点：}

\textbf{有机氯农药（OCPs）：}化学性质极其稳定，高残留农药，脂溶性强，主要蓄积于脂肪组织，生物富集作用强。急性毒性为神经系统和肝脏毒性；慢性毒性影响神经系统、肝脏、血液系统、生殖系统，具有雌激素活性，可经胎盘和乳汁传递，有致畸性。已禁用。

\textbf{有机磷农药（OPPs）：}目前使用范围最广、用量最大的农药。怕光、怕热、怕水，易酶解，半衰期短，残留量低。急性毒性为抑制胆碱酯酶活性，中毒表现为神经功能紊乱，也可出现迟发性神经毒性。慢性毒性影响神经系统、血液系统、视觉。

\textbf{氨基甲酸酯类农药：}毒性低，易被分解，残留量低。毒性为胆碱酯酶抑制剂，但抑制强度较弱且可逆，故毒性比OPPs小，无迟发性神经毒性。可使染色体断裂，有致癌、致畸和致突变作用。

\textbf{拟除虫菊酯类农药：}毒性低，半衰期短，残留量低，但使用后易产生抗药性。急性毒性为神经系统毒性，可引起肌肉痉挛；慢性毒性有致癌性。皮肤接触有刺激和致敏作用，有致突变性和致畸性。
\end{tipbox}

% =====================================================================
%  第七节 $N$-亚硝基化合物污染
% =====================================================================
\section{$N$-亚硝基化合物污染}

\subsection{前体物与合成条件}

\begin{keybox}
\textbf{$N$-亚硝基化合物的前体物*：}
\begin{enumerate}[leftmargin=*]
    \item \textbf{硝酸盐（NO$_3^-$）和亚硝酸盐（NO$_2^-$）：}广泛存在于自然环境（土壤、水体）和食品中；
    \item \textbf{胺类物质：}广泛存在于环境和食品中（蛋白质分解产物）。
\end{enumerate}

\textbf{$N$-亚硝基化合物的体内合成条件*：}
\begin{itemize}[leftmargin=*]
    \item \textbf{pH $<$ 3的酸性环境}是该反应的最适条件。
    \item \textbf{胃是体内合成$N$-亚硝基化合物的最主要场所！}胃液pH约1$\sim$2，且硝酸盐、亚硝酸盐和胺类可随食物同时进入胃内。胃内还有Cl$^-$、SCN$^-$等催化剂。
\end{itemize}
\end{keybox}

\subsection{致癌作用特点}

\begin{keybox}
\textbf{$N$-亚硝基化合物的致癌作用特点*：}
\begin{enumerate}[leftmargin=*]
    \item \textbf{多动物种属诱癌：}对多种实验动物均有致癌性，至今未发现任何动物种属对$N$-亚硝基化合物的致癌作用有抵抗力；
    \item \textbf{多器官/组织靶向性：}对称性亚硝胺主要诱发肝癌，不对称性亚硝胺主要诱发食管癌，可通过血脑屏障和胎盘屏障，引起中枢神经系统肿瘤和经胎盘致癌；
    \item \textbf{多途径均诱癌：}经口、吸入、皮下和肌肉注射、皮肤接触等多种途径均可诱癌；
    \item \textbf{不同接触剂量均可致癌：}一次大剂量或长期小剂量给药均可致癌，有明显的剂量-效应关系；且各剂量组合有相加作用；
    \item \textbf{可以通过胎盘使子代致癌（经胎盘致癌作用）：}怀孕动物摄入后，可通过胎盘影响胎儿，出生后发生肿瘤。
\end{enumerate}
\end{keybox}

\subsection{预防措施}

\begin{keybox}
\textbf{$N$-亚硝基化合物危害的预防*：}
\begin{enumerate}[leftmargin=*]
    \item 防止食品被微生物污染；
    \item 改进食品加工工艺；
    \item 使用钼肥（Mo肥料）——钼是硝酸还原酶的组分，施钼可降低蔬菜中硝酸盐的积累；
    \item 阻断亚硝化反应——多食新鲜蔬菜和水果，增加维生素C、维生素E、酚类、黄酮类化合物的摄入（VC将NO$_2^-$还原为NO，从而阻止亚硝化反应）；
    \item 制定食品中限量标准并加强监测：水产品中$N$-亚硝基化合物限量4 $\mu$g/kg，肉制品中$N$-亚硝基化合物限量3 $\mu$g/kg。% [补充自往届笔记]
\end{enumerate}
\end{keybox}

% =====================================================================
%  第八节 多环芳烃（PAH）污染
% =====================================================================
\section{多环芳烃（PAH）污染}

\begin{keybox}
\textbf{多环芳烃（Polycyclic Aromatic Hydrocarbons, PAH）：}由两个或两个以上苯环融合而成的化合物。\textbf{苯并(a)芘（Benzo(a)pyrene, B(a)P）}是最重要、最具有代表性的PAH。

\textbf{B(a)P的特性：}
\begin{itemize}[leftmargin=*]
    \item 主要蓄积于脂肪组织；
    \item 致胃癌、致突变；
    \item 预防措施：防止污染；去除毒素（活性炭吸附）；制定食品中限量标准（熏烤肉 $\leq$ 5 $\mu$g/kg，植物油 $\leq$ 10 $\mu$g/kg）。% [补充自往届笔记]
\end{itemize}
\end{keybox}

\subsection{污染来源}

\begin{keybox}
\textbf{B(a)P污染食品的主要途径：}
\begin{enumerate}[leftmargin=*]
    \item \textbf{烘烤、烟熏过程中直接污染——最主要的来源！}熏烟中含有大量PAH，烟熏时间越长、温度越高，B(a)P含量越高。炭火烤肉/鱼，油脂滴落在炭火上，热解聚合生成PAH，随烟雾附在食品表面。
    \item 食品成分在高温下的热解/热聚合：脂肪在 $>$ 400$^\circ$C 可生成PAH。
    \item 农作物从环境中吸收：空气沉降的PAH被作物叶片和果实表面吸附。
    \item 食品加工过程中接触污染：如用沥青铺路晒粮（严重违规）。
\end{enumerate}
\end{keybox}

% =====================================================================
%  第九节 杂环胺（Heterocyclic Amines）污染
% =====================================================================
\section{杂环胺（Heterocyclic Amines, HCA）污染}

\begin{keybox}
\textbf{杂环胺污染食品的来源及影响因素：}
\begin{enumerate}[leftmargin=*]
    \item \textbf{烹调方式：}加热温度和时间——温度越高、时间越长、水分含量越少，杂环胺生成越多；
    \item \textbf{食物成分：}蛋白质含量——蛋白质含量越高，杂环胺生成越多，美拉德反应是前提。
\end{enumerate}

\textbf{致癌性：}诱发肝癌、致突变（形成DNA加合物）。

\textbf{预防措施：}
\begin{enumerate}[leftmargin=*]
    \item 改变不良烹调方式和饮食习惯；
    \item 多吃蔬菜水果——蔬菜、水果中的多酚、类黄酮等植物化学物可抑制杂环胺的致突变作用；
    \item 加强监测。
\end{enumerate}
\end{keybox}

% =====================================================================
%  第十节 氯丙醇、丙烯酰胺
% =====================================================================
\section{氯丙醇、丙烯酰胺}

\subsection{氯丙醇（Chloropropanols）}

\begin{defbox}
\begin{itemize}[leftmargin=*]
    \item \textbf{首次发现：}在酸水解植物蛋白液（HVP）中发现，用于酱油等调味品中。
    \item \textbf{代表性物质：}3-氯-1,2-丙二醇（3-MCPD）含量最高、最常见。
    \item \textbf{毒性：}3-MCPD具肾毒性和雄性生殖毒性。
    \item \textbf{存在形式：}游离态3-MCPD（酸水解HVP）以及脂肪酸酯形式（即3-MCPD酯，广泛存在于精炼油脂中）。
\end{itemize}
\end{defbox}

\subsection{丙烯酰胺（Acrylamide）}

\begin{defbox}
\begin{itemize}[leftmargin=*]
    \item \textbf{生成条件：}富含淀粉的食品在 $>$ 120$^\circ$C 高温加工（油炸、烘烤）过程中，由天冬酰胺（asparagine）和还原糖（葡萄糖、果糖）通过Maillard反应生成。
    \item \textbf{高含量食品：}薯条（French fries）、薯片（potato chips）、方便面、油条、速溶咖啡。土豆制品的丙烯酰胺含量约为谷物制品的\textbf{4倍}。
    \item \textbf{毒性和致癌性：}神经毒性（职业暴露），啮齿动物多器官致癌。IARC Group 2A（很可能对人类致癌）。
    \item \textbf{预防：}改进加工工艺（降低油炸温度和时间、避免过度金黄焦脆），原料选用低还原糖的马铃薯品种。
\end{itemize}
\end{defbox}

% =====================================================================
%  第十一节 食品容器、包装材料的卫生问题
% =====================================================================
\section{食品容器、包装材料的卫生问题}

\subsection{概念、范围与分类}

\begin{keybox}
\textbf{食品容器、包装材料的概念：}指在食品生产、加工、储存、运输、销售过程中，直接或间接接触食品的各种材料、制品及辅助材料。

\textbf{范围：}包括从原料到成品的全过程中所有接触食品的器具、管道、设备和包装物。

\textbf{分类及其基本卫生问题（掌握）：}
\begin{enumerate}[leftmargin=*]
    \item \textbf{塑料制品：}PET（聚对苯二甲酸乙二醇酯）、PC（聚碳酸酯）、PE（聚乙烯）、PP（聚丙烯）、PS（聚苯乙烯）、三聚氰胺-甲醛树脂（密胺）、脲醛树脂。基本卫生问题：单体残留（氯乙烯、苯乙烯、双酚A）、增塑剂迁移（邻苯二甲酸酯类——内分泌干扰物）、热稳定性不足。
    \item \textbf{橡胶制品：}天然橡胶、合成橡胶（奶嘴、密封圈）。基本卫生问题：硫化促进剂、抗氧化剂和填充剂的迁移。
    \item \textbf{纸、竹、木：}传统材料。基本卫生问题：漂白剂残留、荧光增白剂、印刷油墨迁移。
    \item \textbf{金属、搪瓷、陶瓷、玻璃：}基本卫生问题：金属离子（铅、镉、铬、镍、铝）的溶出，釉彩中的铅镉迁移。
\end{enumerate}
\end{keybox}

\subsection{影响食品受容器污染的因素}

\begin{defbox}
\begin{enumerate}[leftmargin=*]
    \item \textbf{接触时间和温度：}接触时间越长、温度越高，迁移量越大。热水或热食会显著加速塑料中有害物质溶出。
    \item \textbf{食品的性质：}酸性食品（如果汁、醋）、含酒精食品和油脂类食品对容器材料中有害物质的萃取能力更强。pH越低、脂肪含量越高，迁移量越大。
    \item \textbf{容器材料的性质：}材料的化学组成、加工工艺、添加剂种类和用量决定其卫生安全性。再生塑料的卫生风险较高。
    \item \textbf{接触面积与食品体积之比：}接触面积越大、食品体积越小，单位食品中迁移物的浓度越高。
    \item \textbf{材料的老化程度：}长期使用后材料表面出现裂纹和磨损，增加有害物质的释放。
\end{enumerate}
\end{defbox}

\subsection{主要卫生问题}

\begin{enumerate}[leftmargin=*]
    \item \textbf{有害物质向食品迁移（Migration）——最核心的卫生问题！}
    \begin{itemize}[leftmargin=*]
        \item 塑料中的增塑剂（如邻苯二甲酸酯类，phthalates——具有雌激素样活性，属于内分泌干扰物）
        \item 塑料中的稳定剂、抗氧化剂
        \item 密胺制品（仿瓷餐具）：游离甲醛毒性
        \item 铝制品：盛装酸性或碱性食品可溶解氧化铝膜，铝溶出（铝摄入过量与阿尔茨海默病的关联虽未最终定论，但已引起广泛关注）
    \end{itemize}
    \item 包装材料的回收和再生问题：回收塑料中的不明污染物转移入食品。
    \item 印刷油墨：油墨中的有害物质（如苯系物）可能迁移至食品接触面。
    \item 纳米包装材料：新型材料的安全性评价体系尚待完善。
\end{enumerate}

\subsection{相关卫生标准与指标}

\begin{defbox}
我国已制定GB 4806系列《食品安全国家标准 食品接触材料及制品》标准体系，涵盖：
\begin{itemize}[leftmargin=*]
    \item 各类食品接触材料的通用安全要求
    \item 特定迁移量（SML）和总迁移量（OML）限量
    \item 食品接触材料中添加剂的使用标准（GB 9685）
    \item 各类材料的具体卫生标准（如GB 4806.1--GB 4806.12）
\end{itemize}
\end{defbox}

% =====================================================================
%  复习题
% =====================================================================
\reviewcontent{
\section{复习题}

\begin{enumerate}[leftmargin=*, itemsep=8pt]
    \item \textbf{【名词解释】}食品污染；食品污染物；水分活性（Aw）；菌落总数；大肠菌群；MPN法；霉菌毒素；细菌菌相；霉菌菌相
    \item \textbf{【名词解释】}食品腐败变质；TVBN；K值；D值；食品保藏
    \item \textbf{【名词解释】}农药；农药残留；MRLs；兽药
    \item \textbf{【简答·重点】}简述食品污染的三大分类（生物性、化学性、物理性）及其具体内容。
    \item \textbf{【简答·重点】}简述水分活性（Aw）的定义公式，并说明Aw与微生物生长的关系（含各类微生物的最低Aw阈值及细菌芽孢形成的特殊要求）。
    \item \textbf{【简答·重点】}简述常见食品细菌（假单胞菌属、微球菌属、芽孢杆菌属、肠杆菌科）及其与腐败变质的关系。
    \item \textbf{【简答·重点】}简述细菌菌相的定义及食品卫生学意义。
    \item \textbf{【简答·重点】}简述菌落总数和大肠菌群的食品卫生学意义，并说明MPN法的适用条件及原理。
    \item \textbf{【简答·重点】}简述粪便污染指示菌应具备的条件。
    \item \textbf{【简答·重点】}简述霉菌产毒的条件（基质、水分、湿度、温度、通风）和特点（四项），以及霉菌毒素中毒的共同特点。
    \item \textbf{【简答·重点】}简述霉菌菌相及其食品卫生学意义。
    \item \textbf{【简答·重点】}简述黄曲霉毒素的化学结构、理化性质、毒性排序、产毒条件、体内代谢过程及毒性表现（急性、慢性、致癌性）。为何说黄曲霉毒素是已知最强的天然致癌物？
    \item \textbf{【简答·重点】}详述黄曲霉毒素预防措施，包括食物防霉和7种去毒方法。
    \item \textbf{【简答·重点】}简述镰刀菌毒素三大类（单端孢霉烯族化合物、玉米赤霉烯酮、伏马菌素）的来源及毒性特点。
    \item \textbf{【简答·重点】}简述食品腐败变质的概念、三大原因及各类营养素的腐败产物特征（蛋白质、脂肪、碳水化合物）。
    \item \textbf{【简答·重点】}详述食品腐败变质的鉴定指标（感官、物理、化学、微生物），含TVBN、三甲胺、组胺、K值、pH、过氧化值/酸价的具体卫生学意义。
    \item \textbf{【简答·重点】}简述腐败变质食品的处理原则（四种情况）。
    \item \textbf{【简答·重点】}详述食品保藏的各类方法和关键参数：化学保藏（盐腌浓度、酸渍pH）、低温保藏（冷藏与冷冻的温度范围、冰晶生成带）、高温杀菌（常压/高压/UHT参数、D值及实例、微波频率）、干燥脱水保藏（水分和Aw范围、褐变类型）、辐照保藏（三种剂量分级及用途）。
    \item \textbf{【简答】}简述农药的五大分类方式及各类农药的毒性特点。
    \item \textbf{【简答·重点】}简述$N$-亚硝基化合物的前体物、体内合成条件、致癌作用特点及至少四项预防措施。
    \item \textbf{【简答】}简述B(a)P污染食品的主要途径、蓄积特点和防控措施。
    \item \textbf{【简答·重点】}简述杂环胺污染食品的来源及影响因素、致癌性及预防措施。
    \item \textbf{【简答】}简述丙烯酰胺的生成条件、高含量食品及预防措施。
    \item \textbf{【简答】}简述氯丙醇的来源、代表物质及毒性。
    \item \textbf{【简答·重点】}简述食品容器包装材料的概念、分类及各类的基本卫生问题。
    \item \textbf{【简答】}简述影响食品受容器污染的因素（五项）。
    \item \textbf{【综合论述】}试从食品污染的三大来源出发，综合分析我国食品安全面临的主要化学性污染问题及其综合防控策略。
    \item \textbf{【综合论述】}某地区疑似发生黄曲霉毒素污染事件，请设计一份食品安全调查与防控方案。
\end{enumerate}

\subsection*{参考答案要点}

\begin{enumerate}[leftmargin=*, itemsep=6pt]
    \item \textbf{食品污染：}食品中含有对健康造成危害或影响食用及商品价值的生物性、化学性及物理性物质。\textbf{食品污染物：}在种植、养殖、生产、加工、储存、运输、销售、烹饪、食用等环节进入食品的有害因素。\textbf{Aw：}$\text{Aw}=P/P_0$，同温度下食品水分蒸汽压与纯水蒸汽压之比（0$\sim$1），反映可被微生物利用的实际含水量。\textbf{菌落总数：}单位质量/容积/表面积检样在规定条件下培养后生成的细菌菌落总数（CFU）。\textbf{大肠菌群：}能发酵乳糖、产酸产气、需氧/兼性厌氧的革兰阴性无芽孢杆菌，包含埃希菌属、柠檬酸杆菌属、肠杆菌属和克雷伯菌属。\textbf{MPN法：}最可能数法，基于泊松分布的间接计数法，用于大肠菌群数量低时的统计。\textbf{霉菌毒素：}霉菌污染食品后产生的有毒代谢产物。\textbf{细菌菌相：}共存于食品中的细菌种类及其相对数量构成，优势菌种决定腐败特征。\textbf{霉菌菌相：}食品中污染的真菌种类及其相对数量构成，田野霉占优势示粮食新鲜，毛霉/根霉/木霉检出示已开始霉变。

    \item \textbf{食品腐败变质：}在以微生物为主的各种环境因素作用下食品成分被分解、破坏，失去或降低食用价值的一切变化。\textbf{TVBN：}弱碱条件下水蒸气蒸馏出的总氮量，鱼肉类蛋白腐败鉴定的化学指标。\textbf{K值：}ATP分解的HxR$+$Hx占ATP系列分解产物的百分比，$K<20\%$新鲜，$K\geq40\%$开始腐败。\textbf{D值：}在一定温度下杀死食品中某种细菌90\%的时间（min），如金黄色葡萄球菌D$_{63}$=7，肉毒梭菌芽孢D$_{121}$=0.1$\sim$0.23。\textbf{食品保藏：}为防止食品腐败变质，对微生物进行杀灭或抑制其生长繁殖，通过改变温度、水分、氢离子浓度、渗透压及其他抑菌杀菌措施，延长食品可食用期限。

    \item \textbf{农药：}用于预防控制农林病、虫、草、鼠等有害生物及调节植物昆虫生长的化学合成或生物源物质及其制剂。\textbf{农药残留：}因使用农药在食品中出现的农药衍生物，包括转化物、代谢物、反应产物及杂质。\textbf{MRLs：}按GAP使用农药后在食品中允许的农药残留最高浓度（mg/kg）。\textbf{兽药：}用于预防、治疗、诊断动物疾病或调节动物生理机能的物质，包括药物饲料添加剂。

    \item \textbf{三大分类：}（1）生物性污染：微生物（细菌及毒素、霉菌及毒素、病毒）、寄生虫和昆虫，细菌及其毒素最常见最重要；（2）化学性污染：农药兽药残留、工业三废中汞镉铅砷等有害金属及有机污染物、食品接触材料迁移物、滥用食品添加剂、加工储存中产生的$N$-亚硝基化合物/PAH/杂环胺/丙烯酰胺、掺假制假物质（苏丹红、三聚氰胺）；（3）物理性污染：杂物污染（草籽、灰尘等）和放射性污染（核爆炸、核废料、核事故，水-水生生物-动植物-人转移）。

    \item \textbf{Aw与微生物生长关系：}Aw $= P/P_0$。（1）每种微生物有最低Aw要求，低于则生长受抑；（2）Aw $<$ 0.60绝大多数微生物无法生长；（3）细菌需Aw $>$ 0.9，酵母菌 $>$ 0.87，霉菌 $>$ 0.8；（4）同属不同种Aw要求不同；（5）细菌形成芽孢时需要比生长时更高的Aw值。

    \item \textbf{常见食品细菌：}假单胞菌属（需氧无芽孢杆菌，冷藏食品腐败主要细菌）；微球菌属（需氧无芽孢杆菌，常见于水、蔬菜腐败）；微球菌和葡萄球菌属（革兰阳性，产过氧化氢酶、色素，分解碳水化合物）；芽孢杆菌属与梭状芽孢杆菌属（革兰阳性芽孢杆菌，罐头食品常见腐败菌）；肠杆菌科（需氧无芽孢杆菌，水产、肉、蛋腐败有关，含大肠埃希菌——粪便污染指示菌，变形杆菌属——腐败菌代表）。

    \item \textbf{细菌菌相：}食品中细菌种类及相对数量构成，优势菌种（属、株）决定腐败特征。卫生学意义：（1）影响腐败变质特征，使食用价值降低甚至不能食用；（2）影响腐败变质速度与程度，腐败过程中可出现霉菌繁殖并产毒。

    \item \textbf{菌落总数：}（1）食品细菌污染程度及卫生状态标志；（2）预测食品耐储性——细菌越多腐败越快。\textbf{大肠菌群：}（1）粪便污染指示菌——典型大肠杆菌=近期污染，非典型=陈旧污染；（2）肠道致病菌指示菌——同源同存活时间；（3）检测加工者/设备/存储条件污染情况。\textbf{MPN法：}大肠菌群数量低时采用，基于泊松分布的间接计数法，GB/T 4789.3-2016，单位MPN/g或MPN/mL；数量高时用CFU平板计数法。

    \item \textbf{粪便污染指示菌条件：}（1）仅来自肠道；（2）肠道中数量多、易检出；（3）外环境能生存一定期间；（4）检验方法敏感简易。\textbf{肠球菌（粪链球菌）：}反映低温类食品粪便污染，对低温抵抗力比大肠菌群强。

    \item \textbf{霉菌产毒条件：}基质（营养丰富，天然食品更易繁殖）、水分（粮食17\%$\sim$18\%最适，Aw $<$ 0.7一般不能生长）、湿度（RH $<$ 70\%不能产毒）、温度（最适25$\sim$30$^\circ$C，0$^\circ$C以下或30$^\circ$C以上产毒能力下降）、通风（大部分需氧，毛霉等可厌氧，耐受高浓度CO$_2$）。\textbf{产毒特点：}（1）只限于少数产毒菌株，不同菌株产毒能力不同；（2）产毒能力可变可易变；（3）产毒无严格专一性；（4）产毒需要一定时间。\textbf{霉菌毒素中毒特点：}毒素稳定、经天然食品中毒、实质器官受损为主、无传染性和免疫性、季节性和地区性。

    \item \textbf{霉菌菌相：}田野霉（气连孢霉等）占优势=粮食新鲜；曲霉和青霉检出=储藏条件可能不良；毛霉、根霉、木霉检出=已开始霉变。

    \item \textbf{AF化学结构：}二呋喃氧杂萘邻酮。\textbf{毒性排序：}B1 $>$ M1 $>$ G1 $>$ B2 $>$ M2，监测以AFB1为指标。\textbf{理化性质：}耐热（268$\sim$269$^\circ$C才分解），脂溶性，碱性条件内酯环开环破坏，紫外线下产生荧光。\textbf{产毒条件：}最适温度25$\sim$33$^\circ$C，最适Aw 0.93$\sim$0.98。\textbf{代谢：}AFB1在肝脏CYP450催化下：AFB1 $\rightarrow$ AFM1（羟化，乳汁中）；AFB1 $\rightarrow$ AFQ1（羟化，解毒）；AFQ1 $\rightarrow$ AFH1（还原，解毒）；AFB1 $\rightarrow$ AFB1-8,9-环氧化物（终致癌物，与DNA形成加合物）。\textbf{毒性：}急性中毒性肝炎（黄疸、发热、腹水），慢性肝损伤（发育障碍、肝功降低），致癌性——肝癌，诱癌强度是二甲基亚硝胺75倍，多物种多器官致癌，灵长类可诱癌，与HBV协同致癌，IARC Group 1。AFM1在乳及乳制品中稳定，一般家庭加热及加工过程均对其水平无显著影响。

    \item \textbf{预防：}（1）食物防霉——控制水分在安全水分以下，低温保藏，通风；（2）去毒7法：挑选霉粒法、碾轧加工法、加水搓洗法、植物油加碱去毒（NaOH水解内酯环形成香豆素钠盐，碱炼后水洗）、物理吸附法（活性白陶土/活性炭，液体食品效果好，固体食品效果不明显）、氨气处理法（18kg压力/72$\sim$82$^\circ$C，去除98\%$\sim$100\%）、日光/紫外线照射法；（3）制定限量标准。

    \item \textbf{单端孢霉烯族化合物：}T-2毒素（ATA病原物），化学性质稳定，强细胞毒性/免疫抑制/致畸/潜在致癌性，污染麦类、玉米。\textbf{玉米赤霉烯酮（F-2毒素）：}雌激素样作用，生殖系统毒性，污染玉米、小麦、大麦。\textbf{伏马菌素：}FB1为主，干扰鞘脂类代谢，神经毒性（马脑白质软化症），肾脏毒性，促癌物诱发原发性肝癌，污染玉米。

    \item \textbf{腐败变质概念：}微生物为主各因素作用下，食品成分被分解破坏，失去或降低食用价值。\textbf{三大原因：}（1）微生物——细菌（分解蛋白质）、霉菌和酵母菌（分解碳水化合物、脂肪）；（2）食物本身——酶作用、营养成分和水分（高蛋白=蛋白质腐败，高碳水=产酸发酵，Aw越小越不易腐败）、pH（酸性/碱性不利细菌）；（3）环境——温度、湿度、氧气、光照。

    \textbf{蛋白质腐败：}蛋白质 $\rightarrow$ 多肽 $\rightarrow$ 氨基酸 $\rightarrow$ 脱羧/脱氨 $\rightarrow$ 酪胺、组胺、尸胺（恶臭）+ 硫醇（恶臭）+ 色胺、甲基色胺（粪臭）。\textbf{脂肪酸败：}水解+氧化，中性脂肪 $\rightarrow$ 甘油+脂肪酸 $\rightarrow$ 醛、酮、酸，不饱和脂肪酸氧化形成过氧化物再分解为醛酮酸，不饱和脂肪酸含量越高越易氧化，产生"哈喇味"。酸败早期即可出现酸度增高。鉴定：油脂气味、过氧化值上升（最早期指标）、酸价和羰基价阳性。\textbf{碳水化合物分解：}$\rightarrow$ 单糖、醇、醛、酮、CO$_2$和水，酸度升高，产生酸、醇、醛气味。

    \item \textbf{鉴定指标：}（1）感官——最敏感，色/香/味/形（不良气味、变色、软化、发黏）；（2）物理——油脂比重、折光率、黏度；（3）化学——TVBN（鱼肉类蛋白腐败鉴定，与腐败程度有明确对应关系）；三甲胺（鱼虾水产品新鲜程度）；组胺（$>$200 mg/100g可致过敏性食物中毒）；K值（$<$20\%新鲜，$\geq$40\%开始腐败，鉴定鱼类早期腐败）；pH（V字形变动）；过氧化值（脂肪酸败最早期指标）+ 酸价（脂肪酸败程度）；（4）微生物——菌落总数、大肠菌群、细菌菌相。

    \item \textbf{处理原则：}严重腐败变质——销毁或工业用；轻度腐败——加工处理降至安全限量以下，限期食用；已变质——禁止食用；局部变质——挑选去除变质部分，利用完好部位。

    \item \textbf{化学保藏：}盐腌——食盐浓度达10\%抑制大多数细菌，糖含量达60\%$\sim$65\%即可；酸渍——pH 4.5以下抑制微生物，用于各种蔬菜（泡菜、酸菜）；防腐剂——苯甲酸、山梨酸；抗氧化剂——防止油脂酸败。

    \textbf{低温保藏：}冷藏——$-1\sim10^\circ$C，低温抑制微生物繁殖和酶活性；冷冻——$\leq-18^\circ$C，几乎全部微生物不繁殖；急冻缓化有利于保持食品品质。冰晶生成带：$-1\sim-5^\circ$C（水结冰率85\%），$-8\sim-120^\circ$C（水结冰率90\%），速冻迅速通过冰晶生成带，冰晶数量多、体积小，不损伤组织细胞。

    \textbf{高温杀菌：}常压——100$^\circ$C以下，巴氏杀菌（低温63$^\circ$C/30min，高温短时72$^\circ$C/15s），处理牛乳/啤酒/酱油/葡萄酒；高压——100$\sim$121$^\circ$C（0.2 MPa），肉类制品/罐头，杀灭繁殖型和芽孢型；UHT——封闭系统加热至120$^\circ$C以上数秒后迅速冷却；微波——915 MHz（含水量高/体积大食品）、2450 MHz（含水量低食品），快速节能对品质影响小。

    \textbf{D值：}杀死90\%细菌所需时间（min）。金黄色葡萄球菌D$_{63}$=7，肉毒梭菌芽孢D$_{121}$=0.1$\sim$0.23。

    \textbf{干燥脱水：}水分降至15\%以下或Aw 0.00$\sim$0.60，抑制各种微生物生长。方法：日晒、阴干、喷雾干燥、减压蒸发、冷冻干燥。酶性褐变——酚酶催化酚类形成醌。非酶褐变/Maillard反应——氨基与羰基反应，Aw 0.6最适于褐变，加热与给氧促进褐变，SO$_2$可阻止褐变。

    \textbf{辐照保藏：}辐照灭菌——高剂量10$\sim$15/50 kGy，杀灭所有微生物达到商业无菌；辐照消毒——中等剂量5$\sim$10 kGy，杀灭无芽孢致病菌（沙门菌等），3天内不发虫霉变；辐照防腐——低剂量1$\sim$5 kGy，杀灭腐败菌，延长后熟期及保藏期（抑制发芽、杀虫）。

    \item \textbf{五大分类方式：}按用途（杀虫/杀鼠/杀菌剂等）、按化学组成及结构（有机磷/有机氯/氨基甲酸酯/拟除虫菊酯/有机砷/有机汞等）、按成分和来源（无机/有机/生物农药）、按急性毒性大小（剧毒/高毒/中等毒/低毒）、按残留特性（高残留/中等残留/低残留）。

    \item \textbf{$N$-亚硝基化合物：}前体物为硝酸盐/亚硝酸盐+胺类。合成条件为pH $<$ 3的酸性环境（胃是主要合成场所，胃液pH 1$\sim$2，Cl$^-$/SCN$^-$催化）。致癌特点：（1）多物种无抵抗；（2）多器官靶向（对称亚硝胺=肝癌，不对称=食管癌，可过血脑/胎盘屏障）；（3）多途径诱癌（经口/吸入/注射/皮肤）；（4）单次大剂量或长期小剂量均可致癌，有剂量-效应关系，各剂量组合相加；（5）经胎盘致子代肿瘤。预防：防微生物污染、改进加工工艺、施钼肥、多食新鲜蔬果（VC/VE/酚类/黄酮阻断亚硝化）、制定限量标准（水产品$\leq$4$\mu$g/kg，肉制品$\leq$3$\mu$g/kg）。

    \item \textbf{B(a)P：}主要蓄积于脂肪组织，致胃癌/致突变。污染途径：烘烤/烟熏直接污染（最主要，脂肪滴落炭火热解聚合）、高温热解热聚合（脂肪$>$400$^\circ$C）、农作物从环境吸收、加工接触。防控：防止污染、活性炭吸附去除、制定限量标准（熏烤肉$\leq$5$\mu$g/kg，植物油$\leq$10$\mu$g/kg）。

    \item \textbf{杂环胺：}来源及影响因素——烹调方式（温度越高/时间越长/水分越少生成越多）、食物成分（蛋白质越高生成越多，美拉德反应是前提）。致癌性——肝癌、致突变（DNA加合物）。预防——改变不良烹调方式和饮食习惯、多吃蔬菜水果、加强监测。

    \item \textbf{丙烯酰胺：}富淀粉食品 $>$ 120$^\circ$C高温加工时天冬酰胺+还原糖经Maillard反应生成。高含量食品：薯条、薯片、方便面、油条、速溶咖啡（土豆制品约为谷物制品4倍）。神经毒性，IARC Group 2A。预防：降低油炸温度和时间、选用低还原糖马铃薯品种。

    \item \textbf{氯丙醇：}首次发现于酸水解植物蛋白液（HVP），3-MCPD含量最高最常见，具肾毒性和雄性生殖毒性。游离态3-MCPD（酸水解HVP）及脂肪酸酯形式（3-MCPD酯，广泛存在于精炼油脂中）。

    \item \textbf{食品容器包装材料：}概念——直接或间接接触食品的各种材料、制品及辅助材料。分类：（1）塑料——单体残留（氯乙烯/苯乙烯/双酚A）、增塑剂迁移（邻苯二甲酸酯类，内分泌干扰物）；（2）橡胶——硫化促进剂/抗氧化剂/填充剂迁移；（3）纸竹木——漂白剂残留/荧光增白剂/印刷油墨迁移；（4）金属搪瓷陶瓷玻璃——铅镉铬镍铝溶出，釉彩铅镉迁移。

    \item \textbf{影响迁移的因素：}（1）接触时间和温度——越长/越高迁移量越大；（2）食品性质——酸性/含酒精/油脂类萃取能力更强，pH越低/脂肪越高迁移越大；（3）容器材料性质——化学组成/加工工艺/添加剂，再生塑料风险高；（4）接触面积与食品体积之比——面积越大/体积越小，迁移浓度越高；（5）材料老化程度——裂纹和磨损增加释放。

    \item \textbf{综合防控策略：}（要点）三大污染来源——生物性（微生物：强化HACCP/GMP、冷链管理；寄生虫：加强检验检疫）；化学性（农药：登记制度+安全间隔期+MRLs+高效低毒新品种；有害金属：消除工业污染源是根本+GB 2762限量+金属间拮抗营养策略；$N$-亚硝基化合物：防微生物污染+改进加工工艺+施钼肥+多食新鲜蔬果+限量标准；PAH：改进烟熏工艺+活性炭吸附+限量标准；杂环胺：合理烹调降低温度+多吃蔬菜水果）；物理性（放射性：核事故应急+食品监测；杂物：严格GMP）。综合防控：完善法规标准体系、建立全程可追溯制度、加强风险监测和风险评估、开展食品安全健康教育。

    \item \textbf{黄曲霉毒素污染调查与防控方案：}（要点）调查：采样检测可疑食品（花生/玉米/谷物/植物油等）中AFB1含量，现场调查储藏条件（温度/湿度/Aw），了解人员中毒情况（肝功能检测/流行病学调查）。防控：立即封存/销毁超标食品；查找污染源头（田间/储藏/加工环节）；防霉措施（控制水分至安全水分以下+低温保藏+通风）；已污染食品去毒（挑选霉粒/碾轧加工/碾轧加工/碱炼/吸附/氨气处理/紫外线照射）；严格执行限量标准（GB 2761）；加强健康教育（HBV疫苗接种+AF暴露风险告知，因AF与HBV有协同致癌作用）。
\end{enumerate}

}
\newpage
